\chapter{ДИЗАЙН И ПРОВЕДЕНИЕ ЭКСПЕРИМЕНТОВ}

% -------------------------------------------------------
\section{Обзор экспериментального фреймворка}
% -------------------------------------------------------

\subsection{Три парадигмы как единый спектр}

Настоящее исследование объединяет три классические психологические парадигмы в едином экспериментальном фреймворке, обеспечивающем сопоставимость результатов на уровне моделей, метрик и условий доступа. Выбор парадигм обусловлен концептуальной логикой: они образуют прогрессирующий спектр сложности рассуждения и типа когнитивного предубеждения.

\textit{Эксперимент 1 (ошибка конъюнкции)} тестирует вероятностное суждение в одношаговом формате: модели предъявляется описание персонажа и два варианта ответа, один из которых нормативно правильный. Ошибка возникает в момент выбора и не требует многошагового вывода.

\textit{Эксперимент 2 (задача выбора карточек Уэйсона)} тестирует дедуктивное рассуждение: модель должна применить правило modus tollens к условному утверждению и выбрать карточки, способные его фальсифицировать. Задача является одношаговой, однако требует явного логического анализа структуры правила.

\textit{Эксперимент 3 (задача открытия правила 2-4-6)} тестирует индуктивное рассуждение в многоходовом интерактивном формате: модель должна итеративно формировать и пересматривать гипотезы на основе бинарной обратной связи. Предубеждение проявляется не в единственном выборе, а в стратегии тестирования на протяжении всего диалога.

Все три эксперимента проводились на одном и том же наборе из девяти моделей, что позволяет сравнивать профили предубеждений между парадигмами на уровне отдельных систем.

\subsection{Общие принципы дизайна}

Три принципа обеспечивают методологическую связность фреймворка.

\textit{Контроль контаминации.} В каждом эксперименте используются оригинальные датасеты, не воспроизводящие текст классических задач или их известных адаптаций. Детали контроля контаминации описаны в соответствующих подразделах.

\textit{Воспроизводимость.} Все модели запрашивались через единый API-интерфейс (OpenRouter). Сырые ответы сохранялись отдельно от производных метрик. Случаи нарушения формата фиксировались и обрабатывались через retry-логику.

\textit{Сопоставимость метрик.} Наряду со специфическими для каждого эксперимента показателями вводится единая метрика Severity of Bias, обеспечивающая сквозное сравнение интенсивности предубеждений между парадигмами.

% -------------------------------------------------------
\section{Эксперимент 1: Ошибка конъюнкции (Задача Линды)}
% -------------------------------------------------------

\subsection{Обоснование выбора парадигмы}

В качестве экспериментальной парадигмы для исследования вероятностных предубеждений была выбрана задача ошибки конъюнкции, впервые описанная Тверски и Канеманом~\cite{tversky1983extensional}. В оригинальном эксперименте испытуемым предъявлялось биографическое описание персонажа по имени Линда --- молодой женщины, охарактеризованной как умная, откровенная, обеспокоенная вопросами социальной справедливости. После прочтения испытуемые оценивали относительную вероятность двух утверждений: (a)~Линда является банковским кассиром и (b)~Линда является банковским кассиром и активной участницей феминистского движения. Около 85--90\% испытуемых выбирали вариант (b), нарушая тем самым базовое правило теории вероятностей $P(A \cap B) \leq P(A)$.

Задача была выбрана по трём причинам. Во-первых, она обладает чётко определённым нормативным ответом, что позволяет однозначно классифицировать ответы модели. Во-вторых, механизм ошибки --- нарративная правдоподобность конъюнктивного элемента --- поддаётся прямой экспериментальной манипуляции через замену этого элемента. В-третьих, задача широко использовалась в предшествующих исследованиях~\cite{binz2023using, suri2024large, wang2024linda, jiang2024peek, macmillan2024irrationality}, что обеспечивает возможность сравнения результатов.

\subsection{Проблема контаминации и подходы к её решению}

Задача Линды является одной из наиболее цитируемых задач в психологической литературе и с высокой вероятностью присутствует в обучающих данных большинства современных LLM~\cite{zhou2023benchmark}. Для решения проблемы контаминации был сформирован оригинальный датасет с использованием фиксированных демографических шаблонов и операции строковой замены. Этот подход обеспечивает, с одной стороны, структурное единообразие всех виньеток (исключая конфаундеры, связанные с различиями в нарративном стиле), с другой --- их отсутствие в обучающих корпусах в конкретных использованных комбинациях.

\subsection{Конструкция датасета: биографические виньетки}

Датасет состоит из 110 оригинальных биографических виньеток, составленных вручную. Каждая виньета описывает персонажа с выраженными личностными характеристиками, создающими сильный нарративный образ, и содержит плейсхолдеры \texttt{\{NAME\}}, \texttt{\{RACE\}} и \texttt{\{AGE\}}, заменяемые в ходе подстановки. Пример виньетки:

\begin{quote}
\textit{``\{NAME\} is a \{RACE\} \{AGE\}-year-old public defender who takes on cases others refuse. They grew up in a low-income neighborhood and were the first in their family to attend college. They are known for their passionate speeches about systemic inequality and regularly volunteer at a community legal aid clinic.''}
\end{quote}

К каждой виньете прилагаются пять элементов:
\begin{itemize}
  \item \textbf{T} --- простая опция, содержащая профессию персонажа. Профессии сэмплировались из базы статистики занятости USDL/OEWS~\cite{usdl2024oews}, насчитывающей около 830 профессиональных категорий;
  \item \textbf{F} --- коррелированное хобби, тематически релевантное биографии, но не воспроизводящее дословно ни одного слова из виньетки;
  \item \textbf{F\_uncorrelated} --- некоррелированное хобби, семантически не связанное с биографией персонажа;
  \item \textbf{T\_and\_F1} --- конъюнктивная опция T + F;
  \item \textbf{T\_and\_F2} --- конъюнктивная опция T + F\_uncorrelated.
\end{itemize}

Правило $P(A \cap B) \leq P(A)$ однозначно делает T более вероятной опцией в обоих случаях. Логическая структура задачи идентична в correlated- и uncorrelated-условиях; единственной варьируемой переменной является степень нарративной согласованности второго элемента конъюнкции с биографией персонажа.

\subsection{Демографические подстановки и факториальный дизайн}

Для исследования эффекта демографических токенов использовалась система подстановки имён и возрастов через операцию string replacement. Имена подбирались так, чтобы однозначно ассоциироваться с расовой группой согласно данным о частотности имён в США:

\begin{itemize}
  \item \textbf{White:} Emma, James
  \item \textbf{Black:} Jamal, Aisha
  \item \textbf{Hispanic:} Carlos, Maria
  \item \textbf{Asian:} Wei, Priya
\end{itemize}

Возрасты: 24, 38, 55 лет. Первоначально эксперимент проводился на полной факториальной сетке: 8 имён $\times$ 3 возраста = 24 варианта подстановки на виньетку, что давало более 6000 пар испытаний для первых моделей. Однако анализ результатов не выявил значимых различий по возрасту и proxy-gender: delta\_bias по возрастным группам варьировался в диапазоне $-0{,}005 \div +0{,}003$, групповые тесты МакНемара не прошли порог значимости после коррекции Бенджамини-Хохберга~\cite{benjamini1995controlling}. В связи с этим для последующих моделей сетка была сокращена до расово-значимых подстановок (одна пара имён на расовую группу без разбивки по возрасту), что давало 520 пар испытаний на модель.

\subsection{Протокол эксперимента и промпт-структура}

Для каждой подставленной виньетки проводились два независимых испытания в отдельных сессиях без переноса контекста:

\textbf{Условие A (correlated):} модели предъявлялись опции T и T\_and\_F1. Правильный ответ --- выбор T.

\textbf{Условие B (uncorrelated):} модели предъявлялись опции T и T\_and\_F2. Правильный ответ --- выбор T.

Промпт имел следующую структуру:

\begin{quote}
\textit{Read the following description and answer the question.\\
{[}vignette{]}\\
Which is more probable?\\
(a) {[}T{]}\\
(b) {[}T\_and\_F1 or T\_and\_F2{]}\\
First state your choice: only the letter (a) or (b).\\
Then on a new line state your confidence that the chosen option is more probable, on a scale from 0 to 100, where 100 = absolutely certain.}
\end{quote}

Все модели запрашивались с параметром temperature~=~0. Сырые ответы сохранялись отдельно от распарсенных метрик; случаи нарушения формата логировались и обрабатывались через retry-логику.

\subsection{Метрики: бинарная точность, тест МакНемара, Severity of Bias}

\textbf{Бинарная метрика и тест МакНемара.} Основным инструментом статистического вывода служил точный тест МакНемара для согласованных пар бинарных наблюдений~\cite{mcnemar1947note}. Для каждой пары испытаний (A, B) одной подставленной виньетки строилась таблица 2$\times$2 с ячейками $n_{11}$ (правильно в обоих условиях), $n_{12}$ (правильно в A, ошибочно в B), $n_{21}$ (ошибочно в A, правильно в B), $n_{22}$ (ошибочно в обоих условиях). Нулевая гипотеза предполагает маргинальную однородность: $\pi_{12} = \pi_{21}$. Основная гипотеза предсказывает $n_{21} \gg n_{12}$: модели ошибаются именно тогда, когда хобби семантически связано с биографией персонажа. При $n^* = n_{12} + n_{21} > 10$ использовалась нормальная аппроксимация со статистикой $z_0 = (n_{21} - n_{12}) / \sqrt{n_{21} + n_{12}}$.

\textbf{Severity of Bias.} В дополнение к бинарной классификации ошибок вводится непрерывная метрика силы предубеждения, основанная на самооценке уверенности модели. Для каждого испытания:
\begin{itemize}
  \item если модель выбирала конъюнкцию (b): $\text{bias\_score} = \text{confidence} / 100$;
  \item если модель выбирала T (a): $\text{bias\_score} = (100 - \text{confidence}) / 100$.
\end{itemize}

Шкала от 0 (минимальное предубеждение, высокая уверенность в правильном ответе) до 1 (максимальное предубеждение, высокая уверенность в ошибочном ответе). Ключевой агрегированной метрикой является:
\begin{equation}
  \text{delta\_bias} = \overline{\text{bias}}_{\text{correlated}} - \overline{\text{bias}}_{\text{uncorrelated}},
  \label{eq:delta_bias}
\end{equation}
где $\overline{\text{bias}}_{\text{correlated}}$ и $\overline{\text{bias}}_{\text{uncorrelated}}$ --- средние значения bias\_score в соответствующих условиях. Это значение отражает чистый эффект нарративной связности, очищенный от фонового уровня неуверенности модели. Статистическая значимость delta\_bias проверялась парным t-тестом и критерием Вилкоксона.

\subsection{Результаты и анализ}

Полные результаты эксперимента по всем девяти моделям представлены в таблице~\ref{tab:linda_results}.

\begin{table}[h!]
\caption{Результаты эксперимента 1 по всем моделям}
\label{tab:linda_results}
\centering
\begin{tabular}{|l|c|c|c|c|c|}
\hline
\textbf{Модель} & \textbf{n\_pairs} & \textbf{$n_{21}$} & \textbf{McNemar $p$} & \textbf{delta\_bias} & \textbf{Знак} \\
\hline
qwen3-next-80b   & 520  & 316 & $1{,}50 \cdot 10^{-95}$ & $+0{,}5648$ & + \\
gemma-3-27b-it   & 520  & 305 & $3{,}07 \cdot 10^{-92}$ & $+0{,}4975$ & + \\
gpt-5.2          & 520  & 209 & $2{,}43 \cdot 10^{-63}$ & $+0{,}3746$ & + \\
grok-4.1-fast    & 520  & 119 & $3{,}01 \cdot 10^{-36}$ & $+0{,}2485$ & + \\
deepseek-v3.2    & 520  & 97  & $1{,}26 \cdot 10^{-29}$ & $+0{,}1996$ & + \\
gemma-3-12b-it   & 520  & 14  & $0{,}000122$            & $+0{,}0370$ & + \\
glm-4.7-flash    & 3120 & 26  & $2{,}98 \cdot 10^{-8}$  & $+0{,}0279$ & + \\
claude-sonnet-4.6 & 520 & 1   & $1{,}0$                 & $+0{,}1049$ & +* \\
qwen3.5-397b-a17b & 3120 & 5  & $0{,}0625$              & $-0{,}0052$ & - \\
\hline
\end{tabular}
\end{table}

\textit{Примечание:} $n_{21}$ --- число пар, в которых ошибка совершена только в correlated-условии; знак «+*» для \texttt{claude-sonnet-4.6} означает значимый delta\_bias при бинарно незначимом тесте МакНемара. В столбце McNemar $p$ для \texttt{qwen3.5-397b-a17b} указано $p = 0{,}0625$, что не проходит порог $\alpha = 0{,}05$.

Гипотеза H1 об универсальности токен-bias получила \textit{частичное} подтверждение. Из девяти моделей у восьми выполняется $n_{21} > n_{12} = 0$ и delta\_bias $> 0$, что свидетельствует о систематической склонности ошибаться именно в correlated-условии. Однако для двух моделей результат требует оговорок.

\textbf{claude-sonnet-4.6} демонстрирует феномен «потолка бинарной точности»: за весь прогон из 520 пар допущена единственная бинарная ошибка ($n_{21} = 1$), поэтому тест МакНемара не достигает значимости ($p = 1{,}0$). Вместе с тем метрика delta\_bias фиксирует устойчивый сдвиг: $\text{delta\_bias} = +0{,}1049$ при $p_{\text{t-test}} = 2{,}59 \cdot 10^{-147}$, $p_{\text{Wilcoxon}} = 9{,}60 \cdot 10^{-83}$. Модель систематически демонстрирует сниженную уверенность в правильном ответе именно в correlated-условии, не совершая при этом классификационных ошибок. Это подтверждает, что токен-предубеждение существует на спектре и не сводится к дихотомии «ошибается / не ошибается».

\textbf{qwen3.5-397b-a17b} является единственной моделью, для которой не обнаружено свидетельств токен-bias: $n_{21} = 5$, $n_{12} = 0$, $p_{\text{McNemar}} = 0{,}0625$, delta\_bias $= -0{,}0052$ (отрицательный знак). Статистически значимые $p$-значения t-теста и критерия Вилкоксона ($p = 5{,}27 \cdot 10^{-10}$ и $p = 6{,}27 \cdot 10^{-44}$ соответственно) объясняются не направленным bias, а микросдвигом в калибровке уверенности при практически нулевой вариативности ошибок: модель почти всегда отвечает правильно с высокой уверенностью, и статистически значимое различие между условиями отражает лишь тонкую разницу в распределениях confidence, а не выраженный токен-bias.

\textbf{Внутрисемейный парадокс масштаба} является одним из наиболее неожиданных результатов эксперимента. Внутри семейства Gemma более крупная модель проявляет несравнимо более высокий уровень предубеждения: у \texttt{gemma-3-27b-it} $n_{21} = 305$ и delta\_bias $= 0{,}4975$ против $n_{21} = 14$ и delta\_bias $= 0{,}0370$ у \texttt{gemma-3-12b-it}. Это прямо противоречит гипотезе H2 о монотонном снижении предубеждений с ростом размера модели. Наиболее правдоподобное объяснение: более крупная модель лучше распознаёт нарративную согласованность биографии и конъюнктивного элемента, что парадоксально усиливает эвристику репрезентативности. Меньшая модель, хуже улавливая «смысл» биографии, реже попадает в эту ловушку.

\textbf{Демографический анализ.} По модели \texttt{glm-4.7-flash} зафиксирован расовый эффект: delta\_bias для Asian ($0{,}0318$), Hispanic ($0{,}0314$) и Black ($0{,}0248$) значимо выше нуля после коррекции BH-FDR, тогда как для White отдельный тест МакНемара незначим. По остальным моделям --- \texttt{qwen3.5-397b-a17b}, \texttt{gemma-3-12b-it} --- расовых различий не выявлено. Возраст и proxy-gender (кодировка через имя) не дали значимых эффектов ни у одной модели, что согласуется с решением сократить факториальную сетку для последующих моделей.

\subsection{Выводы по эксперименту 1}

Эксперимент 1 установил: (1)~токен-bias в форме ошибки конъюнкции выявлен у восьми из девяти моделей, что даёт основания говорить о широко распространённом, но не абсолютно универсальном феномене; (2)~метрика Severity of Bias обнаруживает предубеждение там, где бинарная метрика показывает потолок точности, что подтверждено на примере \texttt{claude-sonnet-4.6}; (3)~эффект масштабирования не является монотонным: внутри семейства Gemma более крупная модель демонстрирует многократно более высокий уровень предубеждения; (4)~расовый демографический эффект зафиксирован у \texttt{glm-4.7-flash}, однако не является универсальным; (5)~возраст и proxy-gender не влияют на интенсивность токен-bias.

% -------------------------------------------------------
\section{Эксперимент 2: Задача выбора карточек Уэйсона}
% -------------------------------------------------------

\subsection{Обоснование выбора парадигмы}

Задача выбора карточек Уэйсона~\cite{wason1966reasoning} была выбрана как каноническая парадигма для исследования предубеждения подтверждения в дедуктивном рассуждении. В оригинальной версии испытуемым предъявляются четыре карточки, соответствующие случаям P, $\lnot$P, Q и $\lnot$Q, и условное правило «если P, то Q». Нормативно правильный ответ --- выбор карточек P и $\lnot$Q --- позволяет фальсифицировать правило. Менее 10\% людей решают абстрактную версию правильно; большинство выбирают подтверждающие карточки P и Q.

Задача обеспечивает прямую операционализацию предубеждения подтверждения~\cite{nickerson1998confirmation}: выбор карточек P и Q отражает поиск подтверждающих, а не фальсифицирующих свидетельств. Кроме того, задача позволяет исследовать контент-эффекты~\cite{cheng1985pragmatic, cosmides1992cognitive}: известно, что задачи с содержанием социальных контрактов решаются людьми значительно лучше, чем абстрактные версии.

\subsection{Конструкция стимульных датасетов: четыре контентных условия}

Были сконструированы четыре стимульных датасета, каждый включающий 100 уникальных условных правил с соответствующими отображениями на четыре позиции карточек.

\textbf{Формальная логика --- Каноническая.} Правила используют стандартные абстрактные символы из классической литературы по задаче Уэйсона (буквы, числа, геометрические фигуры). Условие включено для сравнения с предшествующими работами при признании риска контаминации~\cite{binz2023using, arkoudas2023gpt4}.

\textbf{Формальная логика --- Нейтральная.} Правила используют оригинальные абстрактные пары, не присутствующие в канонических формулировках (например, «если символ является Стрелкой, то узор является Полосатым»). Это условие служит деконтаминированным абстрактным базисом и непосредственно адресует методологический пробел работ, опирающихся исключительно на канонические стимулы.

\textbf{Конкретные факты.} Правила основаны на верифицируемых реальных свойствах объектов (например, «если животное является дельфином, то оно дышит воздухом»). Условие операционализирует семантические априорные знания без введения нормативных социальных обязательств.

\textbf{Знакомые социальные контракты.} Деонтические правила, отражающие широко известные социальные нормы и институциональные разрешения (например, «если человек управляет автомобилем, то у него должны быть права»). Условие воспроизводит эффект фасилитации, описанный Ченгом и Холиоаком~\cite{cheng1985pragmatic}. Из набора исключены юрисдикционно-специфические правила.

\textbf{Незнакомые фантастические социальные контракты.} Деонтические правила, вписанные в вымышленные миры, не пересекающиеся с известными художественными вселенными (например, «если путешественник пересекает Лунные ворота, он должен нести Серебряное перо»). Условие разработано для диссоциации двух конкурирующих объяснений эффекта социального контракта: знакомость конкретного содержания правила или распознавание деонтической структуры разрешения/обязательства~\cite{cheng1985pragmatic, cosmides1992cognitive}. Насколько известно автору, подобное условие ранее не реализовывалось в исследованиях LLM.

\subsection{Протокол перестановок порядка карточек}

Ключевым методологическим вкладом дизайна является систематическая манипуляция порядком предъявления карточек. Каждый стимул предъявлялся каждой модели в \textbf{пяти различных перестановках} четырёх позиций карточек. Это следует из наблюдения Binz и Schulz~\cite{binz2023using} о том, что единственное изменение перестановки способно обратить ответ GPT-3 на каноническую задачу, и из данных Macmillan-Scott и Musolesi~\cite{macmillan2024irrationality} о существенной нестабильности ответов LLM.

Данный дизайн даёт \textbf{500 испытаний на модель на условие} (100 правил $\times$ 5 перестановок), позволяя оценивать распределения ответов, а не точечные оценки, и разделять правильность и стабильность как независимые свойства поведения модели.

\subsection{Промпт-условия}

Оценивались три формулировки промпта на датасете «Формальная логика --- Нейтральная» и дополнительно на «Знакомых социальных контрактах».

\textbf{Original (Assisted Falsification).} Явно инструктирует модель искать карточки, способные выявить нарушение правила, в сочетании с персоной «эксперта по логике»:

\begin{quote}
\textit{You are a logic expert analyzing conditional rules. Your task is to identify which cards must be turned over to check if a rule is violated. CRITICAL INSTRUCTIONS: Respond ONLY with the selected cards as a comma-separated list. NO explanations, NO reasoning, NO JSON, NO other text.}
\end{quote}

\textbf{Neutral.} Убирает персону эксперта и фрейминг фальсификации, предъявляя задачу без стратегических указаний. Наиболее близко воспроизводит оригинальную человеческую экспериментальную парадигму.

\textbf{Chain-of-Thought (CoT).} Инструктирует модель рассуждать пошагово перед ответом~\cite{wei2022chain}:

\begin{quote}
\textit{You are a logic expert analyzing conditional rules. Think step by step before answering. After reasoning, output your final answer on the last line as: ANSWER: card1, card2.}
\end{quote}

Условие CoT тестирует, раскрывает ли явное элицирование рассуждений латентную компетентность у моделей, недостаточно показывающих себя при прямых ответах, и взаимодействует ли CoT дифференцированно с уровнем способностей модели~\cite{kojima2022large}.

Все модели запрашивались с temperature~=~1{,}0 для сохранения естественной вариабельности ответов между перестановками. Примеры не предоставлялись; все условия --- zero-shot.

\subsection{Метрики}

\textbf{Expected Metric Accuracy (EMA)} --- доля испытаний, в которых модель выбирает ровно нормативно правильный набор карточек $\{$P, $\lnot$Q$\}$. Основная мера логической правильности, сопоставимая с метриками точности в работах Dasgupta и соавторов~\cite{dasgupta2024language} и Macmillan-Scott и Musolesi~\cite{macmillan2024irrationality}.

\textbf{Matching Bias Index (MBI)} --- доля испытаний, в которых модель выбирает $\{$P, Q$\}$ --- карточки, явно упомянутые в правиле~\cite{evans1998matching}. Операционализирует предубеждение соответствия.

\textbf{Confirmation Bias Rate (CBR)} --- доля испытаний, в которых модель выбирает только $\{$P$\}$ --- минимальный подтверждающий выбор. Совместно с MBI формирует «семейство предубеждения подтверждения».

\textbf{Consistency Rate} --- доля правил, для которых модель даёт идентичный выбор карточек во всех пяти перестановках. Количественно оценивает стабильность ответов независимо от правильности.

\subsection{Результаты и анализ}

\textbf{Результаты по контентным условиям.} Полные данные по всем пяти датасетам приведены в таблицах~\ref{tab:wst_content1} и~\ref{tab:wst_content2}.

\begin{table}[h!]
\caption{EMA по контентным условиям (Formal Logic Canonical и Neutral)}
\label{tab:wst_content1}
\centering
\begin{tabular}{|l|c|c|c|c|}
\hline
\textbf{Модель} & \textbf{FL Canon.} & \textbf{FL Neutral} & \textbf{Conc. Facts} & \textbf{Fam. Social} \\
\hline
gpt-5.2                   & 0,962 & 0,972 & 0,976 & 0,996 \\
qwen3-next-80b            & 0,834 & 0,976 & 0,980 & 0,946 \\
qwen3.5-397b-a17b         & 0,822 & 0,818 & 0,909 & 0,982 \\
claude-sonnet-4.6         & 0,814 & 0,898 & 0,956 & 0,982 \\
deepseek-v3.2             & 0,764 & 0,880 & 0,883 & 0,922 \\
grok-4.1-fast             & 0,634 & 0,720 & 0,875 & 0,818 \\
gemma-3-27b-it            & 0,716 & 0,578 & 0,792 & 0,904 \\
glm-4.7-flash             & 0,464 & 0,614 & 0,562 & 0,686 \\
gemma-3-12b-it            & 0,332 & 0,106 & 0,380 & 0,546 \\
\hline
\end{tabular}
\end{table}

\begin{table}[h!]
\caption{EMA и Consistency Rate в условии Unfamiliar Fantasy Social Contracts}
\label{tab:wst_content2}
\centering
\begin{tabular}{|l|c|c|c|}
\hline
\textbf{Модель} & \textbf{Fantasy Social EMA} & \textbf{Confirmation-fam.} & \textbf{Consistency} \\
\hline
qwen3.5-397b-a17b  & 0,986 & 0,004 & 0,950 \\
gpt-5.2            & 0,978 & 0,010 & 0,920 \\
claude-sonnet-4.6  & 0,970 & 0,024 & 0,920 \\
qwen3-next-80b     & 0,938 & 0,030 & 0,780 \\
gemma-3-27b-it     & 0,888 & 0,104 & 0,750 \\
deepseek-v3.2      & 0,882 & 0,028 & 0,580 \\
grok-4.1-fast      & 0,818 & 0,082 & 0,440 \\
glm-4.7-flash      & 0,590 & 0,334 & 0,260 \\
gemma-3-12b-it     & 0,436 & 0,240 & 0,170 \\
\hline
\end{tabular}
\end{table}

\textbf{Контент-эффект} является наиболее устойчивым результатом эксперимента и воспроизводится у всех девяти моделей. Переход от абстрактного нейтрального условия к знакомым социальным контрактам даёт прирост EMA во всех без исключения случаях: медианный прирост составляет $+0{,}31$ (от $+0{,}164$ у \texttt{qwen3.5-397b-a17b} до $+0{,}440$ у \texttt{gemma-3-12b-it}). Критически важным является результат условия «Незнакомых фантастических социальных контрактов»: EMA в нём значимо превышает EMA в нейтральном абстрактном условии у восьми из девяти моделей и вплотную приближается к EMA знакомых социальных контрактов. Это поддерживает гипотезу прагматических схем~\cite{cheng1985pragmatic}: деонтическая структура обязательства фасилитирует правильный логический вывод независимо от знакомости содержания.

Исключением является \texttt{gemma-3-27b-it}, у которой EMA в условии Fantasy ($0{,}888$) превышает EMA в Familiar ($0{,}904$) незначительно, но при этом EMA в нейтральном абстрактном условии ($0{,}578$) заметно ниже --- что указывает на особую чувствительность этой модели именно к деонтическому фреймингу.

\textbf{Предубеждение подтверждения} подтверждено для всех моделей в нейтральном промпт-условии (таблица~\ref{tab:wst_prompts}). В Formal Logic Neutral при нейтральном промпте \texttt{gemma-3-12b-it} демонстрирует Confirmation-family $= 0{,}670$ (MBI $= 0{,}536$ + CBR $= 0{,}134$) при EMA $= 0{,}106$ --- модель в 6,4 раза чаще выбирает предубеждённый ответ, чем правильный.

\begin{table}[h!]
\caption{Сравнение промпт-условий на датасете Formal Logic Neutral}
\label{tab:wst_prompts}
\centering
\begin{tabular}{|l|c|c|c|c|c|c|}
\hline
\multirow{2}{*}{\textbf{Модель}} & \multicolumn{3}{c|}{\textbf{Neutral prompt}} & \multicolumn{3}{c|}{\textbf{CoT prompt}} \\
\cline{2-7}
 & EMA & MBI & CBR & EMA & MBI & CBR \\
\hline
gpt-5.2            & 0,884 & 0,056 & 0,002 & 0,976 & 0,006 & 0,006 \\
claude-sonnet-4.6  & 0,790 & 0,172 & 0,002 & 0,970 & 0,000 & 0,000 \\
qwen3-next-80b     & 0,772 & 0,038 & 0,024 & 0,880 & 0,000 & 0,004 \\
qwen3.5-397b-a17b  & 0,764 & 0,032 & 0,000 & 0,974 & 0,000 & 0,004 \\
deepseek-v3.2      & 0,658 & 0,200 & 0,002 & 0,936 & 0,000 & 0,002 \\
grok-4.1-fast      & 0,424 & 0,450 & 0,000 & 0,664 & 0,032 & 0,252 \\
gemma-3-27b-it     & 0,406 & 0,386 & 0,006 & 0,368 & 0,436 & 0,120 \\
glm-4.7-flash      & 0,332 & 0,304 & 0,248 & 0,258 & 0,086 & 0,548 \\
gemma-3-12b-it     & 0,028 & 0,270 & 0,134 & 0,186 & 0,468 & 0,230 \\
\hline
\end{tabular}
\end{table}

\textbf{Эффект CoT-промпта} является одним из наиболее содержательных результатов. Для топ-5 моделей CoT даёт существенный прирост EMA: \texttt{claude-sonnet-4.6} (+0,180), \texttt{qwen3.5-397b-a17b} (+0,210), \texttt{deepseek-v3.2} (+0,278), при этом MBI и CBR у всех пяти моделей падают практически до нуля. Однако для нижней части рейтинга CoT не только не помогает, но и вредит. У \texttt{glm-4.7-flash} EMA снижается с 0,332 до 0,258, а CBR вырастает с 0,248 до 0,548: явное рассуждение систематически приводит эту модель к выбору только карточки $\{$P$\}$. У \texttt{gemma-3-27b-it} CoT не меняет EMA содержательно (0,406 vs 0,368), однако перераспределяет ошибки: MBI растёт с 0,386 до 0,436, то есть явное рассуждение усиливает предубеждение соответствия. Это указывает на то, что CoT не является универсальным средством коррекции предубеждений и может давать обратный эффект у моделей с недостаточной базовой логической компетентностью~\cite{kojima2022large}.

\textbf{Стабильность ответов} в Formal Logic Neutral при нейтральном промпте: от 0,110 у \texttt{gemma-3-12b-it} до 0,910 у \texttt{qwen3-next-80b}. Примечательно, что \texttt{qwen3-next-80b} демонстрирует наивысшую стабильность (0,910) при EMA 0,976 --- сочетание, нетипичное для других моделей. \texttt{gemma-3-27b-it} при EMA 0,578 имеет стабильность 0,280 --- низкую относительно точности, что указывает на нестабильное угадывание, а не на систематическое рассуждение.

\textbf{Внутрисемейный парадокс.} В семействе Gemma эффект из эксперимента 1 инвертируется: по EMA в Formal Logic Neutral \texttt{gemma-3-27b-it} ($0{,}578$) значительно превосходит \texttt{gemma-3-12b-it} ($0{,}106$). Таким образом, в вероятностной задаче 27b была более уязвима, чем 12b, а в логической задаче --- менее уязвима. Это указывает на то, что уязвимость к разным типам когнитивных предубеждений не определяется единым фактором и может иметь различную направленность в зависимости от типа рассуждения.

\subsection{Выводы по эксперименту 2}

Эксперимент 2 установил: (1)~предубеждение подтверждения в дедуктивном рассуждении у всех девяти моделей в нейтральном условии; (2)~устойчивый контент-эффект --- социальные контракты решаются значимо лучше абстрактных правил у всех моделей; (3)~поддержку гипотезы прагматических схем Ченга-Холиоака через условие фантастических социальных контрактов, которое даёт EMA, сопоставимое со знакомыми контрактами; (4)~дифференцированный эффект CoT --- сильные модели значительно улучшают результаты, слабые ухудшают; (5)~инверсию внутрисемейного парадокса Gemma по сравнению с экспериментом 1, что свидетельствует о мультидимензиональности когнитивных предубеждений.

% -------------------------------------------------------
\section{Эксперимент 3: Задача открытия правила 2-4-6}
% -------------------------------------------------------

\subsection{Обоснование выбора парадигмы}

Задача открытия правила 2-4-6, введённая Уэйсоном~\cite{wason1960failure}, является классической парадигмой для исследования предубеждения подтверждения в индуктивном рассуждении. В оригинальном эксперименте испытуемым сообщалось, что тройка $(2, 4, 6)$ удовлетворяет скрытому правилу, и предлагалось открыть его, предлагая новые тройки и получая бинарную обратную связь. Канонический результат: испытуемые тестируют тройки, согласующиеся с текущей гипотезой (например, $(4, 8, 10)$ для гипотезы «чётные числа в порядке возрастания»), вместо попыток фальсифицировать её, и не обнаруживают истинное правило («любая возрастающая последовательность»).

В отличие от экспериментов 1 и 2, где предубеждение проявляется в единственном ответе, здесь оно измеряется как свойство \textit{стратегии} на протяжении многоходового диалога.

\subsection{Конструкция датасета правил: пять таксономических категорий}

Был сформирован независимый датасет из \textbf{50 правил}, организованных в пять таксономических категорий:

\begin{itemize}
  \item \textbf{Порядок и сравнение} (9 правил) --- отношения между элементами: $x < y < z$, $x > y \land z > y$;
  \item \textbf{Чётность и делимость} (10 правил) --- свойства модульной арифметики: все чётные, сумма кратна 3;
  \item \textbf{Арифметические отношения} (10 правил) --- аддитивные и мультипликативные структуры: $x + y = z$, $y - x = z - y$;
  \item \textbf{Структурные паттерны} (11 правил) --- позиционные и теоретико-множественные свойства: $\max(x, y, z) = z$, ровно два уникальных значения;
  \item \textbf{Смешанные / составные} (10 правил) --- конъюнкции из предыдущих категорий: $x < y < z \land (x + y + z) \equiv 0 \pmod{2}$.
\end{itemize}

Все правила верифицированы на исполнимость в диапазоне $[1, 99]^3$, отсутствие функциональных дубликатов (проверка через экстенсиональную эквивалентность) и отсутствие тавтологий. Ряд правил намеренно сконструирован как \textit{ловушки специфичности}: правила, являющиеся строго более специфичными, чем истинное правило, но согласующиеся с порождающей тройкой $(2, 4, 6)$. Из 50 правил \textbf{26 удовлетворяются} тройкой $(2, 4, 6)$ и 24 не удовлетворяются.

\subsection{Многоходовой диалоговый протокол}

Эксперимент реализован как многоходовой диалог. На каждом ходу $t$ модель получает историю разговора и должна ответить в строго ограниченном формате:

\begin{quote}
\texttt{H: lambda x, y, z: <hypothesis>}\\
\texttt{T: [x, y, z]}
\end{quote}

где \texttt{H} --- исполняемый Python-лямбда-выражение, представляющее текущую гипотезу о скрытом правиле $R$, а \texttt{T} --- следующая тройка для тестирования. Модель получает только бинарную обратную связь вида \texttt{(a, b, c) -> True} или \texttt{(a, b, c) -> False}. Диалог продолжается не более \texttt{max\_turns} шагов. В отличие от бенчмарка WILT~\cite{banatt2024wilt}, дизайн настоящего исследования требует от модели поддерживать живую, исполняемую гипотезу \textit{на каждом ходу}, что позволяет анализировать полную траекторию гипотез.

\subsection{Промпт-условия: Baseline и Adaptive}

\textbf{Baseline.} Предоставляет только фрейм задачи: существование скрытого правила, порождающую тройку, требуемый формат вывода и жёсткие ограничения (целые числа в $[1, 99]$, только \texttt{H} и \texttt{T} без пояснений).

\textbf{Adaptive.} Расширяет базовый промпт явными метакогнитивными указаниями, основанными на нормативном байесовском пересмотре гипотез: после получения ответа False гипотезу следует сузить; после неожиданного True --- расширить; рекомендуется рассматривать арифметические свойства (суммы, произведения, остатки) как компоненты правила. Манипуляция мотивирована данными Batista и Griffiths~\cite{batista2026rational} о том, что явные инструкции по обратной связи модулируют поведение пересмотра гипотез.

\subsection{Метрики качества гипотез}

\textbf{Ключевой методологической особенностью} является использование \textit{экстенсиональной эквивалентности} для оценки гипотез. Вместо текстового сравнения строк гипотезы оцениваются как множества принятых троек над большой случайной выборкой~\cite{banatt2024wilt}. Формально: для скрытого правила $R$ и гипотезы модели $H_t$ на ходу $t$ сэмплируются $N = 200{,}000$ случайных троек $(x, y, z) \in [1, 99]^3$ и вычисляются булевы векторы $\mathbf{r}$ и $\mathbf{h}_t$. Метрика:
\begin{equation}
  \text{IOU}(H_t, R) = \frac{|\mathbf{h}_t \cap \mathbf{r}|}{|\mathbf{h}_t \cup \mathbf{r}|},
  \label{eq:iou}
\end{equation}
испытание считается \textbf{успешным} при $\text{IOU}(H_T, R) > 0{,}95$.

\textbf{Success Rate} ($\uparrow$) --- доля правил, для которых финальная гипотеза достигает $\text{IOU} > 0{,}95$.

\textbf{Mean AUC-IOU} ($\uparrow$) --- площадь под кривой IOU по ходам, усреднённая по всем правилам:
\begin{equation}
  \text{AUC-IOU} = \frac{1}{T} \sum_{t=1}^{T} \text{IOU}(H_t, R).
  \label{eq:auc_iou}
\end{equation}

\textbf{Confirming Test Rate (CTR)} ($\downarrow$) --- доля ходов, на которых предлагаемая тройка $T_t$ предсказывается текущей гипотезой $H_t$ как True:
\begin{equation}
  \text{CTR} = \frac{1}{T} \sum_{t=1}^{T} \mathbf{1}[H_t(T_t) = \text{True}].
  \label{eq:ctr}
\end{equation}
Напрямую операционализирует предубеждение подтверждения: низкий CTR соответствует фальсификационно-ориентированной стратегии.

\textbf{Hypothesis Change Rate (HCR)} ($\uparrow$) --- доля последовательных пар ходов, на которых модель существенно пересматривает гипотезу:
\begin{equation}
  \text{HCR} = \frac{1}{T-1} \sum_{t=2}^{T} \mathbf{1}[\text{IOU}(H_t, H_{t-1}) < 0{,}9].
  \label{eq:hcr}
\end{equation}
Низкий HCR указывает на «залипание» в начальной гипотезе вне зависимости от полученной обратной связи.

\subsection{Результаты и анализ}

Полные результаты по baseline- и adaptive-условиям представлены в таблице~\ref{tab:246_results}.

\begin{table}[h!]
\caption{Результаты эксперимента 3: baseline и adaptive условия}
\label{tab:246_results}
\centering
\begin{tabular}{|l|c|c|c|c|}
\hline
\textbf{Модель} & \textbf{SR base} & \textbf{SR adap} & \textbf{CTR base} & \textbf{CTR adap} \\
\hline
claude-sonnet-4.6  & 0,28 & 0,28 & 0,812 & 0,787 \\
qwen3.5-397b-a17b  & 0,22 & 0,18 & 0,850 & 0,602 \\
qwen3-next-80b     & 0,20 & 0,24 & 0,976 & 0,719 \\
gpt-5.2            & 0,18 & 0,18 & 0,719 & 0,584 \\
grok-4.1-fast      & 0,14 & 0,08 & 0,629 & 0,672 \\
deepseek-v3.2      & 0,10 & 0,16 & 0,879 & 0,768 \\
gemma-3-12b-it     & 0,06 & 0,04 & 0,747 & 0,718 \\
glm-4.7-flash      & 0,04 & 0,06 & 0,734 & 0,486 \\
gemma-3-27b-it     & 0,04 & 0,10 & 0,710 & 0,842 \\
\hline
gpt-5.2 reasoning  & 0,38 & 0,48 & 0,726 & 0,559 \\
\hline
\end{tabular}
\end{table}

\textbf{Гипотеза H1 подтверждена полностью:} CTR в baseline-условии варьируется от 0,629 у \texttt{grok-4.1-fast} до 0,976 у \texttt{qwen3-next-80b}. Даже наименее «подтверждающая» модель в 63\% случаев тестирует тройки, которые её же текущая гипотеза предсказывает как True. Нормативная стратегия байесовского агента предполагает CTR близкий к нулю.

\textbf{Success Rate низкий у всех non-reasoning моделей.} Лучший результат в baseline среди них --- \texttt{claude-sonnet-4.6} с SR $= 0{,}28$, что означает успешное открытие правила лишь в 14 из 50 случаев. \texttt{gemma-3-27b-it} и \texttt{glm-4.7-flash} открывают правило лишь в двух случаях из 50.

\textbf{Reasoning даёт качественный скачок.} \texttt{gpt-5.2} в режиме reasoning (baseline) достигает SR $= 0{,}38$ --- на 10 процентных пунктов выше лучшей non-reasoning модели. В adaptive-условии SR вырастает до 0{,}48 и CTR снижается до 0{,}559, что является единственным случаем CTR ниже 0{,}60 во всём эксперименте. Таким образом, reasoning + adaptive --- единственная комбинация, приближающаяся к разумной фальсификационной стратегии, хотя и не достигающая нормативного стандарта.

\textbf{Adaptive-промпт нестабилен.} Из девяти non-reasoning моделей у четырёх SR не изменился или снизился (\texttt{gpt-5.2}: без изменений; \texttt{qwen3.5-397b-a17b}: $0{,}22 \to 0{,}18$; \texttt{grok-4.1-fast}: $0{,}14 \to 0{,}08$; \texttt{gemma-3-12b-it}: $0{,}06 \to 0{,}04$). Снижение CTR не гарантирует роста SR.

Наиболее интригующим является \textbf{парадокс \texttt{qwen3.5-397b-a17b}}: adaptive-промпт вызвал крупнейшее снижение CTR среди non-reasoning моделей ($0{,}850 \to 0{,}602$, то есть модель стала значительно активнее фальсифицировать), однако SR упал с 0{,}22 до 0{,}18. Модель начала тестировать опровергающие тройки, но не смогла конвертировать это в правильную гипотезу --- что указывает на разрыв между стратегией тестирования и качеством обновления гипотез.

\textbf{Контринтуитивный эффект у \texttt{gemma-3-27b-it}:} adaptive-промпт, направленный на фальсификацию, \textit{увеличил} CTR ($0{,}710 \to 0{,}842$), однако SR при этом вырос с 0{,}04 до 0{,}10. Это указывает на то, что данная модель находит правильную гипотезу не через фальсифицирующую стратегию, а через какой-то иной механизм, нечувствительный к инструкциям по типу тестов.

\textbf{AUC-IOU} выявляет расхождение между качеством траектории и финальным результатом. У \texttt{gpt-5.2} reasoning при переходе baseline $\to$ adaptive: SR $+0{,}10$, но AUC-IOU $0{,}119 \to 0{,}083$ --- менее гладкая траектория сходимости при лучшем финальном результате. У \texttt{deepseek-v3.2} в baseline AUC-IOU ($0{,}052$) существенно ниже, чем у \texttt{claude-sonnet-4.6} ($0{,}127$), при схожем SR, что указывает на «спринт» к правильному ответу в последних ходах без устойчивого накопления качества гипотезы.

\subsection{Выводы по эксперименту 3}

Эксперимент 3 установил: (1)~предубеждение подтверждения в индуктивном рассуждении универсально --- CTR $\geq 0{,}63$ у всех non-reasoning моделей в baseline; (2)~reasoning даёт качественный скачок, недостижимый через prompt engineering: SR \texttt{gpt-5.2} reasoning = 0{,}48 против максимума 0{,}28 у non-reasoning; (3)~adaptive-промпт нестабилен: снижение CTR не конвертируется в рост SR у половины моделей; (4)~парадокс \texttt{qwen3.5-397b-a17b} --- разрыв между стратегией тестирования и качеством обновления гипотез --- указывает на то, что предубеждение подтверждения имеет как минимум два независимых компонента: выбор тестируемых троек и ревизию гипотез.

% -------------------------------------------------------
\section{Сквозной анализ: три парадигмы, девять моделей}
% -------------------------------------------------------

\subsection{Сравнительная таблица результатов по всем экспериментам}

Таблица~\ref{tab:cross_experiment} представляет агрегированные результаты по всем трём экспериментам для каждой из девяти моделей.

\begin{table}[h!]
\caption{Сводная таблица результатов по трём экспериментам}
\label{tab:cross_experiment}
\centering
\begin{tabular}{|l|c|c|c|c|}
\hline
\textbf{Модель} & \textbf{delta\_bias (Э1)} & \textbf{fallacy rate (Э1)} & \textbf{EMA neutral (Э2)} & \textbf{SR base (Э3)} \\
\hline
qwen3.5-397b-a17b & $-0{,}005$ & $0{,}2\%$  & 0,764 & 0,22 \\
glm-4.7-flash     & $+0{,}028$ & $0{,}8\%$  & 0,332 & 0,04 \\
gemma-3-12b-it    & $+0{,}037$ & $2{,}7\%$  & 0,028 & 0,06 \\
claude-sonnet-4.6 & $+0{,}105$ & $0{,}2\%$  & 0,790 & 0,28 \\
deepseek-v3.2     & $+0{,}200$ & $18{,}7\%$ & 0,658 & 0,10 \\
grok-4.1-fast     & $+0{,}249$ & $22{,}9\%$ & 0,424 & 0,14 \\
gpt-5.2           & $+0{,}375$ & $40{,}2\%$ & 0,884 & 0,18 \\
gemma-3-27b-it    & $+0{,}498$ & $58{,}7\%$ & 0,406 & 0,04 \\
qwen3-next-80b    & $+0{,}565$ & $60{,}8\%$ & 0,772 & 0,20 \\
\hline
gpt-5.2 reasoning & ---        & ---        & ---   & 0,38 \\
\hline
\end{tabular}
\end{table}

\textit{Примечание:} delta\_bias и fallacy rate --- эксперимент 1, чем ниже delta\_bias тем лучше; EMA neutral --- Formal Logic Neutral при нейтральном промпте, эксперимент 2, чем выше тем лучше; SR base --- Success Rate в базовом условии, эксперимент 3, чем выше тем лучше. Таблица отсортирована по delta\_bias. \texttt{gpt-5.2} reasoning выделен отдельно как режим с принципиально иным типом вывода.

\subsection{Проверка гипотез о масштабировании и внутрисемейный парадокс}

Гипотеза H2 о снижении предубеждений с ростом модели получила частичное подтверждение. На уровне общей тенденции зависимость наблюдается: более мощные модели (\texttt{gpt-5.2}, \texttt{claude-sonnet-4.6}) систематически демонстрируют более низкий уровень предубеждений во всех трёх парадигмах. Однако внутри семейств картина неоднородна.

В семействе Gemma направление эффекта противоречит ожидаемому: \texttt{gemma-3-27b-it} демонстрирует несравнимо более высокий токен-bias (delta\_bias $= 0{,}498$, fallacy rate $58{,}7\%$), чем \texttt{gemma-3-12b-it} (delta\_bias $= 0{,}037$, fallacy rate $2{,}7\%$). Это является наиболее ярким проявлением внутрисемейного парадокса масштаба, зафиксированного в эксперименте 1. В семействе Qwen наблюдается противоположный паттерн: \texttt{qwen3-next-80b} при меньшем числе активных параметров демонстрирует существенно более высокий bias (delta\_bias $= 0{,}565$), чем \texttt{qwen3.5-397b-a17b} (delta\_bias $= -0{,}005$). Оба семейства указывают на то, что масштаб сам по себе не является надёжным предиктором устойчивости к когнитивным предубеждениям, и подчёркивают роль post-training alignment.

\subsection{Проверка гипотезы о спектральной природе предубеждений}

Гипотеза H3 подтверждена: уровень предубеждений систематически возрастает от эксперимента 1 к эксперименту 3. В задаче Линды в uncorrelated-условии ошибок нет ни у одной модели (все $n_{12} = 0$); в задаче Уэйсона EMA в нейтральном условии при нейтральном промпте не превышает 0,884 даже у лучшей модели, а у \texttt{gemma-3-12b-it} падает до 0,028; в задаче 2-4-6 CTR $> 0{,}60$ у всех моделей без исключения. Таким образом, индуктивное открытие правил остаётся наиболее уязвимой к предубеждению подтверждения областью для современных LLM.

\subsection{Итоги: Proof-of-Concept фреймворка оценки когнитивных предубеждений}

Совокупность результатов трёх экспериментов позволяет сформулировать следующие итоговые положения.

Во-первых, когнитивные предубеждения в LLM являются устойчивым, воспроизводимым феноменом. Токен-bias зафиксирован у восьми из девяти моделей в эксперименте 1; предубеждение подтверждения в дедуктивном рассуждении --- у всех девяти в эксперименте 2; CTR $\geq 0{,}63$ у всех non-reasoning моделей в эксперименте 3. Гипотеза H1 подтверждена с оговоркой: \texttt{qwen3.5-397b-a17b} не проявляет значимого токен-bias в задаче Линды, однако демонстрирует выраженное предубеждение подтверждения в задачах Уэйсона и 2-4-6 (CTR $= 0{,}850$).

Во-вторых, профили предубеждений по экспериментам не коррелируют монотонно. Наглядный пример: \texttt{gpt-5.2} имеет высокий токен-bias в задаче Линды (fallacy rate $40{,}2\%$, delta\_bias $= 0{,}375$), но лучшую EMA в задаче Уэйсона ($0{,}884$) и наилучший SR в задаче 2-4-6 среди non-reasoning моделей ($0{,}18$, а в reasoning-режиме --- $0{,}38$). Это означает, что уязвимость к разным типам предубеждений не является единым латентным фактором, и ни одна агрегированная метрика не заменяет полного профиля.

В-третьих, масштаб модели не является надёжным предиктором устойчивости к предубеждениям. Внутрисемейный парадокс Gemma проявляется по-разному в разных экспериментах: в задаче Линды \texttt{gemma-3-27b-it} значительно хуже, чем \texttt{gemma-3-12b-it}; в задаче Уэйсона --- значительно лучше (EMA $0{,}578$ vs $0{,}106$); в задаче 2-4-6 обе модели слабы (SR $0{,}04$ и $0{,}06$). Это подчёркивает, что post-training alignment оказывает неоднородное воздействие на разные типы рассуждения.

В-четвёртых, reasoning является единственным вмешательством, дающим качественный скачок: SR \texttt{gpt-5.2} reasoning + adaptive ($0{,}48$) вдвое превышает лучший non-reasoning результат ($0{,}28$). Prompt engineering в форме CoT улучшает результаты только у моделей с достаточной базовой компетентностью и может давать обратный эффект у слабых моделей.